\documentclass[11pt]{article}
\usepackage{lmodern}
\usepackage[T1]{fontenc}
\usepackage[utf8]{inputenc}
\usepackage[margin=28mm]{geometry}
\usepackage{amsmath,amssymb}
\usepackage{booktabs}
\usepackage{xcolor}
\usepackage{fancyhdr}
\usepackage{microtype}
\usepackage{enumitem}
\usepackage{titlesec}
\usepackage[hidelinks,
  pdftitle={EFC Pre-Registered Prediction: EG Statistic from SO x Euclid},
  pdfauthor={Morten Magnusson}]{hyperref}

\definecolor{eblue}{RGB}{0,70,140}
\titleformat{\section}{\large\bfseries\color{eblue}}{}{0em}{}
\titleformat{\subsection}{\normalsize\bfseries\color{eblue}}{}{0em}{}

\pagestyle{fancy}
\fancyhf{}
\rfoot{\small Magnusson --- EFC SO$\times$Euclid $E_G$ Pre-Registration --- Page \thepage}
\renewcommand{\headrulewidth}{0pt}

\setlength{\parskip}{4pt}
\setlength{\parindent}{0pt}

\begin{document}

% Title block
\begin{center}
{\LARGE\bfseries\color{eblue}
Pre-Registered EFC Prediction:\\[4pt]
$E_G$ Statistic from SO\,$\times$\,Euclid Cross-Correlation}

\vspace{8pt}
{\large Morten Magnusson}\\[2pt]
{\small Symbiose Research, Sandnes, Norway\\
ORCID: \href{https://orcid.org/0009-0002-4860-5095}{0009-0002-4860-5095}\\
April 2026 \quad CC-BY-4.0}

\vspace{6pt}
\textbf{Pre-registration date:} 15 April 2026\quad
\textbf{Data expected:} SO\,$\times$\,Euclid DR2 ($\sim$2027--2028)
\end{center}

\vspace{6pt}
\begin{quote}
\textbf{Abstract.}
Energy-Flow Cosmology (EFC) predicts a specific enhancement of the
gravitational slip statistic $E_G$ relative to $\Lambda$CDM, arising
from the simultaneous action of $\mu < 1$ (growth suppression) and
$\Sigma > 1$ (lensing enhancement) derived self-consistently from a single variational
action. We pre-register the quantitative prediction
$E_G^{\rm EFC}/E_G^{\rm GR} = 1.086 \pm 0.012$, together with frozen
pass/fail criteria, the dataset specification (Simons Observatory
LAT\,$\times$\,Euclid DR2), and explicit falsification conditions.
All numerical values are derived from independently archived,
timestamped sources. This document constitutes a sealed prediction
prior to the availability of the targeted observational data.
\end{quote}

\newpage
\section{1\quad Scope and Governance}

This document records a frozen, pre-registered prediction for the
$E_G$ gravitational slip statistic, to be tested against the
cross-correlation of Simons Observatory (SO) CMB lensing with
Euclid DR2 galaxy clustering and weak lensing data.

\textbf{Governance rules (frozen at pre-registration):}
\begin{itemize}[itemsep=2pt]
\item No post-hoc adjustment of the predicted value or pass/fail
  window is permitted after this document is deposited.
\item The EFC parameters ($\mu_0 = 0.94$, $\Sigma_0 = 1.05$,
  $\eta = 1.10$) are frozen in the Global Parameter Registry
  (\href{https://doi.org/10.6084/m9.figshare.31876324}{DOI 31876324}).
\item Falsification is binding: if the kill criterion is met, the
  EFC gravitational slip sector is formally abandoned.
\item All results --- including failures --- will be reported in the
  EFC Validation Ledger within one revision cycle.
\end{itemize}

\textbf{What this prediction does \emph{not} claim:}
This is a consistency test, not a discovery claim. A result consistent
with the EFC prediction does not constitute confirmation of EFC;
it constitutes non-rejection of the gravitational-slip sector.

\section{2\quad Physical Basis}

\subsection{2.1\quad The $E_G$ Observable}

The $E_G$ statistic combines CMB lensing, galaxy--galaxy lensing,
and redshift-space distortions (RSD) to probe the ratio of the
lensing potential to the velocity growth rate:

\begin{equation}
E_G(\ell) = \frac{C_\ell^{\kappa g}}{\beta \, C_\ell^{gg}},
\qquad \beta = \frac{f}{b},
\label{eq:EG}
\end{equation}

where $C_\ell^{\kappa g}$ is the CMB-lensing$\times$galaxy
cross-power spectrum, $C_\ell^{gg}$ the galaxy auto-power,
$f$ the linear growth rate, and $b$ the galaxy bias.
In general relativity, $E_G^{\rm GR} = \Omega_m(z)/f(z)$.
In practice, the observable corresponds to a survey-weighted
estimator and may differ from this idealised expression; all
predictions refer to the pipeline-level definition used in
SO\,$\times$\,Euclid analyses.
The reference quantity $E_G^{\rm GR}$ is evaluated within a
standard $\Lambda$CDM background as used in survey pipelines.
The EFC prediction is therefore defined relative to this
operational baseline, not as a claim about the fundamental
validity of the $\Lambda$CDM background itself.
Both $E_G^{\rm obs}$ and $E_G^{\rm GR}$ are defined within
the same survey pipeline and estimator framework to avoid
inconsistencies arising from mixed definitions.
For a modified gravity theory with $\mu \ne 1$ and $\Sigma \ne 1$:

\begin{equation}
E_G^{\rm MG} = E_G^{\rm GR} \times \frac{\Sigma}{\mu \, f / f_{\rm GR}},
\label{eq:EGMG}
\end{equation}

making $E_G$ the \emph{orthogonal probe} to weak lensing alone:
it is sensitive to the ratio $\Sigma/\mu$, not to each separately.
This relation holds in the quasi-static, sub-horizon regime relevant
for the adopted $\ell$-range, and follows from the modified Poisson
and lensing equations in that limit; see Bellini \& Sawicki (2014)
for the general EFT mapping.

\subsection{2.2\quad EFC Prediction Derivation}

The EFC Relativistic Action
(\href{https://doi.org/10.6084/m9.figshare.31876324}{DOI 31876324})
derives from a single variational principle:
\begin{equation}
S_{\rm EFC} = \int d^4x\,\sqrt{-g}\,\Bigl[
\tfrac{1}{2}M_{\rm Pl}^2 F(\phi)\,R
- \tfrac{1}{2}K(\rho)\,g^{\mu\nu}\partial_\mu\phi\,\partial_\nu\phi
- V(\phi) - \lambda\bigl(\Box\phi - \Gamma(\rho)\bigr)
\Bigr].
\label{eq:action}
\end{equation}

This action produces simultaneously:
\begin{align}
\mu &\approx 0.94 \quad (\text{growth suppression via entropy stiffness}),\\
\Sigma &\approx 1.05 \quad (\text{lensing enhancement via flow anisotropy}),\\
\eta &\approx 1.10 \quad (\text{gravitational slip, non-zero}).
\end{align}

From the Case~B eigenmode analysis
(\href{https://doi.org/10.6084/m9.figshare.31990194}{DOI 31990194}),
the $E_G$ ratio at $k_c \approx 0.05\,h\,{\rm Mpc}^{-1}$ and
$z \approx 0.5$ evaluates to:

\begin{equation}
\boxed{
\frac{E_G^{\rm EFC}}{E_G^{\rm GR}} = 1.086 \pm 0.012
}
\label{eq:prediction}
\end{equation}

where the uncertainty covers the survival valley
$\mu \in [0.93, 0.96]$, $\Sigma \in [1.03, 1.07]$.
The Bellini--Sawicki mapping
(\href{https://doi.org/10.6084/m9.figshare.32011407}{DOI 32011407})
confirms $\alpha_T = 0$ exactly and derives
$\alpha_M \propto S(a)$, $\alpha_B \propto dS/d\ln a$,
establishing EFC compatibility with the Euclid EFT pipeline.

\subsection{2.3\quad Why $E_G$ is the Key Discriminant}

\begin{table}[!ht]
\centering
\caption{Discriminating power of $E_G$ across competing models.}
\vspace{4pt}
\begin{tabular}{@{}lcccc@{}}
\toprule
Model & $\mu < 1$ & $\Sigma > 1$ & $\eta \ne 1$ &
  $E_G^{\rm obs}/E_G^{\rm GR} > 1$ \\
\midrule
EFC                 & \checkmark & \checkmark & \checkmark & \checkmark \\
Massive neutrinos   & \checkmark &            &            &            \\
$f(R)$ gravity      &            &            &            &            \\
Horndeski (general) & possible   & possible   &            & possible   \\
$w_0 w_a$CDM        &            &            &            &            \\
\bottomrule
\end{tabular}
\label{tab:discriminant}
\end{table}

No currently established model listed in Table~\ref{tab:discriminant} produces
$\mu < 1$ \emph{and} $\Sigma > 1$ \emph{and} $\eta \ne 1$
simultaneously from a single action. A detection of
$E_G^{\rm obs}/E_G^{\rm GR} \approx 1.09$ would provide strong
evidence in favour of EFC-class theories over all alternatives
in the table.

\newpage
\section{3\quad Dataset Specification}

\begin{table}[!ht]
\centering
\caption{Pre-registered dataset specification.}
\vspace{4pt}
\begin{tabular}{@{}ll@{}}
\toprule
Element & Specification \\
\midrule
CMB lensing & Simons Observatory LAT $\kappa$ reconstruction,
  $\ell \in [100, 2000]$ \\
Galaxy sample & Euclid DR2 photometric catalogue,
  $z \in [0.3, 0.8]$, 5 tomographic bins \\
$\ell$-range & $\ell \in [50, 500]$ (linear regime; $k < 0.2\,h\,{\rm Mpc}^{-1}$) \\
Pivot scale & $k_c = 0.05\,h\,{\rm Mpc}^{-1}$, $z_{\rm eff} = 0.5$ \\
Bias treatment & Survey-provided NLA + linear bias marginalisation \\
Covariance & Survey-provided covariance matrices; no EFC nuisance parameters \\
Expected date & Euclid DR2 + SO cross-correlation $\sim$2027--2028 \\
\bottomrule
\end{tabular}
\label{tab:dataset}
\end{table}

\textbf{Note on tomographic binning:} Euclid tomographic window
functions are expected to reduce the intrinsic signal amplitude
by $\sim 30$--$40\%$. The pass/fail window in
Section~\ref{sec:criteria} is defined for the \emph{tomography-convolved}
observable; no correction is applied post-hoc.

\section{4\quad Pass/Fail Criteria}
\label{sec:criteria}

\begin{table}[!ht]
\centering
\caption{Pre-registered pass/fail criteria (frozen, no post-hoc revision).}
\vspace{4pt}
\begin{tabular}{@{}lp{5.0cm}p{4.3cm}p{2.0cm}@{}}
\toprule
Criterion & Threshold & Consequence & Regime \\
\midrule
\textbf{KC-EG-PASS} &
  $E_G^{\rm obs}/E_G^{\rm GR} \in [1.03,\,1.14]$ at $> 1\sigma$ &
  Consistent with EFC prediction & FA4 supported \\[4pt]
\textbf{KC-EG-FAIL} &
  $E_G^{\rm obs}/E_G^{\rm GR} = 1.00 \pm 0.01$ at $> 3\sigma$ &
  EFC gravitational-slip sector \textbf{falsified} & FA4 triggered \\[4pt]
\textbf{KC-EG-INC.} &
  Result between pass and fail windows &
  No action; await Euclid DR3 & --- \\
\bottomrule
\end{tabular}
\label{tab:criteria}
\end{table}

\textbf{Kill criterion FA4} (from the Stage-IV Roadmap,
\href{https://doi.org/10.6084/m9.figshare.31970904}{DOI 31970904}):
if $|\eta - 1| < 0.01$ is established at $> 3\sigma$ from the
SO\,$\times$\,Euclid EG measurement, the EFC perturbation sector
is formally falsified and the result is permanently documented
in the Validation Ledger.

\textbf{Note on pass window width:} The pass window $[1.03, 1.14]$
is intentionally broader than the $1\sigma$ prediction interval
to account for survey window functions, tomographic convolution,
and estimator-level effects that are not modelled post-hoc.
This is a pre-registered design choice, not a soft target.

\section{5\quad Regime Declaration and Systematics}

\textbf{RCMP regime declaration.}
Following the Regime-Consistent Measurement Principle
(\href{https://doi.org/10.6084/m9.figshare.31222900}{DOI 31222900}),
the three components of $E_G$ are classified by regime:
CMB lensing (L1 linear, large-scale);
galaxy clustering (L1/L2 boundary, bias-dependent);
RSD (L1 growth sector).
This declaration makes explicit a consistency condition that is
typically implicit in standard analyses.
Cross-regime mixing is declared here, not corrected for post-hoc.
All predictions are evaluated at the intersection of these regimes
at $k_c = 0.05\,h\,{\rm Mpc}^{-1}$, $z_{\rm eff} = 0.5$,
where all three components remain in the linear regime.
The prediction is stated at a representative pivot scale rather than
as a full functional form, to maintain strict pre-registration
without introducing additional model freedom.

\begin{table}[!ht]
\centering
\caption{Known systematics and EFC response.}
\vspace{4pt}
\begin{tabular}{@{}lp{4cm}p{5cm}@{}}
\toprule
Systematic & Effect on $E_G$ & EFC robustness \\
\midrule
Intrinsic alignments & Shifts $C_\ell^{\kappa g}$ amplitude &
  Marginalised; NLA model per survey spec \\[3pt]
Photo-$z$ errors & Dilutes tomographic signal &
  Scale cuts at $k < 0.2\,h\,{\rm Mpc}^{-1}$ applied \\[3pt]
Galaxy bias & Enters via $\beta = f/b$ &
  Cancels in $E_G$ ratio to leading order \\[3pt]
Magnification bias & Sub-percent at $\ell < 500$ &
  Below signal threshold \\[3pt]
Baryonic feedback & Confined to $k > 0.5\,h\,{\rm Mpc}^{-1}$ &
  Outside pre-registered $\ell$-range \\[3pt]
Massive neutrinos & Suppress $\mu$ but $\eta = 1$ &
  Discriminated by $\eta \ne 1$ measurement \\[3pt]
RSD modelling & Affects $\beta = f/b$ estimation &
  Mitigated via survey pipeline; cancels partially in $E_G$ ratio \\[3pt]
CMB lensing reconstruction & Bias in $C_\ell^{\kappa g}$ from map filtering &
  Mitigated via standard SO pipeline null tests and simulations \\
\bottomrule
\end{tabular}
\label{tab:systematics}
\end{table}

\section{6\quad Reproducibility}

\begin{sloppypar}
\begin{itemize}[itemsep=2pt]
\item \textbf{RCMP} (regime-consistent measurement principle):
  \href{https://doi.org/10.6084/m9.figshare.31222900}{DOI 31222900}
\item \textbf{EFC Relativistic Action} (derivation of $\mu$, $\Sigma$, $\eta$):
  \href{https://doi.org/10.6084/m9.figshare.31876324}{DOI 31876324}
\item \textbf{Case B eigenmode} ($E_G^{\rm EFC}/E_G^{\rm GR} = 1.086$):
  \href{https://doi.org/10.6084/m9.figshare.31990194}{DOI 31990194}
\item \textbf{Bellini--Sawicki mapping} (hi\_class compatibility):
  \href{https://doi.org/10.6084/m9.figshare.32011407}{DOI 32011407}
\item \textbf{Stage-IV Roadmap} (KC4 kill criterion):
  \href{https://doi.org/10.6084/m9.figshare.31970904}{DOI 31970904}
\item \textbf{EFC Pre-Registered Predictions} (sealed parameter set):
  \href{https://doi.org/10.6084/m9.figshare.32010399}{DOI 32010399}
\item \textbf{GitHub repository}:
  \href{https://github.com/supertedai/EFC}{github.com/supertedai/EFC}
\item \textbf{Validation Ledger}:
  \href{https://supertedai.github.io/EFC/public/EFC_Validation_Ledger.html}{EFC Validation Ledger}
\end{itemize}
\end{sloppypar}

All code and data are version-controlled. The git commit hash at the
time of analysis will be recorded in the Validation Ledger entry.
No code modification is permitted after unblinding.

\vspace{8pt}
\hrule
\vspace{6pt}
\begin{sloppypar}
\small
\textbf{Final statement.}
This document constitutes a complete, frozen pre-registration for the
EFC $E_G$ prediction prior to the availability of SO\,$\times$\,Euclid
DR2 data. The predicted value $E_G^{\rm EFC}/E_G^{\rm GR} = 1.086 \pm 0.012$,
the pass/fail window, the dataset specification, and the falsification
conditions are not subject to modification. Pre-registration date:
15 April 2026.

\vspace{4pt}
\copyright\ 2026 Morten Magnusson. Licensed under CC-BY-4.0.
\end{sloppypar}

\end{document}
